% Clase 9 — la caricatura de la orientacion molecular, y su desmentida.
% El docente dibuja las moleculas alineadas y ACTO SEGUIDO aclara que es
% mentira ("ojo, ojo, imagen, caricatura"). La figura tiene que mostrar las dos
% cosas o repite el error que el se preocupa por corregir: de ahi el tercer
% panel, que es el que dice la verdad.
% Los dipolos se dibujan como un par de puntos unidos; el angulo de cada uno se
% fija a mano y no al azar, para que el panel "realista" se vea desordenado pero
% con un sesgo leve y reproducible.
\begin{tikzpicture}[scale=1]

  % ---------- panel 1: sin campo ----------
  \begin{scope}[shift={(0,0)}]
    \draw[figgray!45, line width=0.6pt] (-1.5,-1.5) rectangle (1.5,1.5);
    \foreach \x/\y/\a in {-1.0/1.0/25, -0.1/1.05/200, 0.9/0.95/110,
                          -1.05/0.1/300, -0.05/0.0/70, 1.0/0.15/160,
                          -0.95/-1.0/135, 0.0/-1.05/340, 1.05/-0.95/55}{
      \begin{scope}[shift={(\x,\y)}, rotate=\a]
        \draw[figline, line width=0.7pt] (-0.22,0) -- (0.22,0);
        \node[figneg, minimum size=4pt] at (-0.22,0) {};
        \node[figpos, minimum size=4pt] at (0.22,0) {};
      \end{scope}}
    \node[figlbl] at (0,-1.95) {sin campo};
    \node[figlblaux, align=center] at (0,-2.6) {orientaciones\\ aleatorias};
  \end{scope}

  % ---------- panel 2: la caricatura ----------
  \begin{scope}[shift={(4.4,0)}]
    \draw[figgray!45, line width=0.6pt] (-1.5,-1.5) rectangle (1.5,1.5);
    \foreach \x in {-1.0,0.0,1.0}{
      \foreach \y in {-1.0,0.0,1.0}{
        \begin{scope}[shift={(\x,\y)}, rotate=90]
          \draw[figline, line width=0.7pt] (-0.22,0) -- (0.22,0);
          \node[figneg, minimum size=4pt] at (-0.22,0) {};
          \node[figpos, minimum size=4pt] at (0.22,0) {};
        \end{scope}}}
    \node[figlbl] at (0,-1.95) {la caricatura};
    \node[figlblaux, figred, align=center] at (0,-2.6) {\emph{no} es lo que pasa};
  \end{scope}

  % ---------- panel 3: lo que realmente pasa ----------
  \begin{scope}[shift={(8.8,0)}]
    \draw[figgray!45, line width=0.6pt] (-1.5,-1.5) rectangle (1.5,1.5);
    % mismos angulos que el panel 1, sesgados ~15 grados hacia la vertical
    \foreach \x/\y/\a in {-1.0/1.0/40, -0.1/1.05/195, 0.9/0.95/105,
                          -1.05/0.1/295, -0.05/0.0/85, 1.0/0.15/150,
                          -0.95/-1.0/125, 0.0/-1.05/335, 1.05/-0.95/70}{
      \begin{scope}[shift={(\x,\y)}, rotate=\a]
        \draw[figline, line width=0.7pt] (-0.22,0) -- (0.22,0);
        \node[figneg, minimum size=4pt] at (-0.22,0) {};
        \node[figpos, minimum size=4pt] at (0.22,0) {};
      \end{scope}}
    \node[figlbl] at (0,-1.95) {lo que pasa};
    \node[figlblaux, align=center] at (0,-2.6) {casi igual de desordenadas,\\ apenas sesgadas};
  \end{scope}

  % ---------- campo aplicado, comun a los paneles 2 y 3 ----------
  \foreach \x in {3.2,5.6,7.6,10.0}{
    \draw[figvecaux] (\x,-1.5) -- (\x,1.5);}
  \node[figlbl, figgray] at (6.6,1.85) {$\vec E$ aplicado};

  \node[figlbl, align=center] at (4.4,-3.6)
    {El promedio del panel 1 da cero; el del panel 3, no. \textbf{Ese} promedio no nulo};
  \node[figlbl, align=center] at (4.4,-4.15)
    {es la polarizaci\'on --- no el alineamiento perfecto del panel 2.};

\end{tikzpicture}
